\documentclass[border=10pt]{standalone}
\usepackage{fontspec}
\AtBeginDocument{\hyphenation{zenodo}}% never hyphenate "zenodo" inside a DOI
\setmainfont{Noto Serif}[Ligatures=TeXOff]
\XeTeXgenerateactualtext=1
\newfontfamily\bengalifont{Noto Serif Bengali}[Script=Bengali, NFSSFamily=bnfam]
\usepackage[Latin,Bengali]{ucharclasses}
% group-free font switching (as in header-common.tex): robust inside braces, cells and quotes
\makeatletter
\def\BN@saved{\rmdefault}
\def\BN@bn{bnfam}
\newcommand\BNin{\edef\BN@tmp{\f@family}\ifx\BN@tmp\BN@bn\else\let\BN@saved\BN@tmp\fi\fontfamily{bnfam}\selectfont}
\newcommand\BNout{\fontfamily{\BN@saved}\selectfont}
\makeatother
\setTransitionsFor{Bengali}{\BNin}{\BNout}
% Typographic quotes and dashes are not in a ucharclasses class (class 0), so a Bangla word right
% after “ or ‘ would not trigger the Bangla font. Treat them as Latin so that every switch fires.
\XeTeXcharclass"2018=\XeTeXcharclass`A
\XeTeXcharclass"2019=\XeTeXcharclass`A
\XeTeXcharclass"201C=\XeTeXcharclass`A
\XeTeXcharclass"201D=\XeTeXcharclass`A
\XeTeXcharclass"2013=\XeTeXcharclass`A
\XeTeXcharclass"2014=\XeTeXcharclass`A
\XeTeXcharclass"2026=\XeTeXcharclass`A
\usepackage{newunicodechar}
\newunicodechar{।}{{\bengalifont\char"0964}}
\newunicodechar{॥}{{\bengalifont\char"0965}}

\usepackage[table]{xcolor}
\usepackage{booktabs,tabularx,array}
\definecolor{bapilrule}{HTML}{1F2937}
\definecolor{bapilstripe}{HTML}{F3F4F6}
\renewcommand{\arraystretch}{1.28}
\setlength{\tabcolsep}{5pt}

\begin{document}
\begin{minipage}{17cm}
{\large\bfseries সারণি ৩৬।}\ {\large রেফারেন্স বাস্তবায়নের মডিউলসমূহ}\par\vspace{7pt}
\small
\rowcolors{2}{bapilstripe}{white}
\begin{tabularx}{\linewidth}{@{}>{\hsize=0.819\hsize\raggedright\arraybackslash}X>{\hsize=1.735\hsize\raggedright\arraybackslash}X>{\hsize=0.771\hsize\raggedright\arraybackslash}X>{\hsize=0.675\hsize\raggedright\arraybackslash}X@{}}
\arrayrulecolor{bapilrule}\toprule
\rowcolor{white}\textbf{মডিউল} & \textbf{কী করে} & \textbf{প্রয়োজনীয়তা} & \textbf{পরীক্ষা} \\
\midrule
bapil.text & বাংলার ইউনিকোড প্রমিতকরণ, অঙ্ক একীকরণ, হালকা বাংলা প্রত্যয় অপসারণ, দ্বিভাষিক শব্দকোষ সম্প্রসারণ & FR-01, NFR-01 & test\_text \\
bapil.corpus & Markdown নথি লোড, উদ্ধৃতিযোগ্য অনুচ্ছেদে বিভাজন, উৎস-তথ্য ও SHA-256 চেকসাম সংরক্ষণ & FR-03 & test\_corpus \\
bapil.search & ভাষা-সচেতন টোকেনাইজেশন ও আন্তঃভাষিক প্রশ্ন-সম্প্রসারণসহ অনুচ্ছেদভিত্তিক BM25 সূচি & FR-01 & test\_search \\
bapil.answer & নিষ্কাশনমূলক উত্তর, যার প্রতিটি বাক্যে উদ্ধৃতি থাকে; প্রমাণ-সীমার নিচে বিরত থাকে & FR-02, NFR-08 & test\_answer \\
bapil.risk & সংযুক্তি খ-এর ঝুঁকি-স্তর প্রশ্নমালা ও স্কোরিং নিয়ম; স্তর, কারণ ও আবশ্যিক নিয়ন্ত্রণ প্রদান & FR-07, GR-02, GR-06 & test\_risk \\
bapil.registry & JSON স্কিমায় যাচাইকৃত এআই ব্যবস্থা নিবন্ধন; জনসমক্ষে স্বচ্ছতা নিবন্ধন রপ্তানি & FR-06, FR-09 & test\_registry \\
bapil.evaluate & সংস্করণকৃত প্রশ্নসেটে Recall@k, MRR ও nDCG; JSON প্রতিবেদন & FR-08 & test\_evaluate \\
bapil.audit & যাচাই-সুবিধাসহ হ্যাশ-শৃঙ্খলিত, হস্তক্ষেপ-শনাক্তযোগ্য JSONL নিরীক্ষা লগ & GR-04 & test\_audit \\
bapil.api & Python স্ট্যান্ডার্ড লাইব্রেরিভিত্তিক ন্যূনতম JSON HTTP API (অনুসন্ধান, উত্তর, ঝুঁকি) & FR-10 & test\_api \\
bapil.cli & কমান্ড লাইন: index, search, ask, assess, evaluate, register, verify-log & সব & test\_cli \\
\bottomrule
\end{tabularx}
\par\vspace{3pt}{\scriptsize\color{gray} উৎস: BAPIL টেকনিক্যাল ব্লুপ্রিন্ট, সংস্করণ ১.০ (২০২৬)। CC BY 4.0।}
\end{minipage}
\end{document}
